% Section 4: Gravitational Theory
\section{Gravitational Theory}
\label{sec:gravity}

\subsection{From Riemann to Riemann-Cartan Geometry}

General relativity is based on Riemannian geometry, where the connection is symmetric ($\Gamma^\lambda_{\mu\nu} = \Gamma^\lambda_{\nu\mu}$) and torsion vanishes. Our theory extends this to Riemann-Cartan geometry, where torsion is non-zero and carries physical significance.

In Riemann-Cartan geometry:
\begin{itemize}
    \item Metric compatibility: $\nabla_\lambda g_{\mu\nu} = 0$
    \item Non-symmetric connection: $\Gamma^\lambda_{\mu\nu} - \Gamma^\lambda_{\nu\mu} = T^\lambda{}_{\mu\nu}$
    \item Curvature with torsion: $R^\lambda{}_{\mu\nu\rho}(\Gamma) = \partial_\nu\Gamma^\lambda_{\mu\rho} - \partial_\rho\Gamma^\lambda_{\mu\nu} + \Gamma^\lambda_{\sigma\nu}\Gamma^\sigma_{\mu\rho} - \Gamma^\lambda_{\sigma\rho}\Gamma^\sigma_{\mu\nu}$
\end{itemize}

The connection decomposes as:
\begin{equation}
    \Gamma^\lambda_{\mu\nu} = \{^\lambda_{\mu\nu}\} + K^\lambda{}_{\mu\nu}
\end{equation}
where $\{^\lambda_{\mu\nu}\}$ is the Christoffel symbol and $K^\lambda{}_{\mu\nu}$ is the contortion tensor:
\begin{equation}
    K^\lambda{}_{\mu\nu} = \frac{1}{2}(T^\lambda{}_{\mu\nu} + T_{\mu\nu}{}^\lambda + T_{\nu\mu}{}^\lambda)
\end{equation}

\subsection{Einstein-Cartan Field Equations}

The gravitational field equations in our framework are the Einstein-Cartan equations:

\begin{align}
    G_{\mu\nu}(\{g\}) &= 8\pi G T^{(m)}_{\mu\nu} \\
    T^{\lambda\mu\nu} + 2g^{\lambda[\mu}T^{\nu]} &= 8\pi G \tau_0^2 S^{\lambda\mu\nu}
\end{align}

where $S^{\lambda\mu\nu}$ is the spin angular momentum tensor of matter.

In the low-energy limit ($\tau_0 \to 0$), these reduce to standard Einstein equations with vanishing torsion.

\subsection{Six Gravitational Wave Polarizations}

Unlike general relativity's two polarization modes (plus and cross), our theory predicts six distinct polarization modes:

\begin{enumerate}
    \item \textbf{$h_+$ (Plus)}: $e_{xx} = -e_{yy} = 1$ — standard GR mode
    \item \textbf{$h_\times$ (Cross)}: $e_{xy} = e_{yx} = 1$ — standard GR mode
    \item \textbf{$h_x$ (Vector-x)}: $e_{xz} = e_{zx} = 1$ — UFT prediction
    \item \textbf{$h_y$ (Vector-y)}: $e_{yz} = e_{zy} = 1$ — UFT prediction
    \item \textbf{$h_b$ (Breathing)}: $e_{xx} = e_{yy} = 1$ — UFT prediction
    \item \textbf{$h_l$ (Longitudinal)}: $e_{zz} = 1$ — UFT prediction
\end{enumerate}

The relative amplitudes are:
\begin{equation}
    h_+ : h_\times : h_x : h_y : h_b : h_l = 1 : 1 : \frac{\tau_0}{2} : \frac{\tau_0}{2} : \frac{\tau_0^2}{3} : \frac{\tau_0^2}{5}
\end{equation}

For $\tau_0 = 10^{-6}$:
\begin{equation}
    h_x/h_+ \sim 5 \times 10^{-7}, \quad h_b/h_+ \sim 3 \times 10^{-13}
\end{equation}

\subsection{Propagation Speeds and Dispersion}

The six polarization modes propagate at slightly different speeds:
\begin{align}
    v_{+,\times} &= c \\
    v_{x,y} &= c(1 - 10^{-1}\tau_0^2) \approx c(1 - 10^{-13}) \\
    v_b &= c(1 - 2 \times 10^{-1}\tau_0^2) \approx c(1 - 2 \times 10^{-13}) \\
    v_l &= c(1 - 3 \times 10^{-1}\tau_0^2) \approx c(1 - 3 \times 10^{-13})
\end{align}

Over cosmological distances ($d \sim 1$ Gpc), this leads to arrival time differences:
\begin{equation}
    \Delta t \sim \frac{d}{c} \times 10^{-13} \sim 10^{-5} \text{ s}
\end{equation}

This is detectable with LISA's timing precision ($\sim 10^{-9}$ s).

\subsection{Weak Field Limit}

In the weak field limit ($g_{\mu\nu} = \eta_{\mu\nu} + h_{\mu\nu}$), the linearized field equations are:
\begin{equation}
    \Box \bar{h}_{\mu\nu} = -16\pi G T_{\mu\nu}
\end{equation}

where $\bar{h}_{\mu\nu} = h_{\mu\nu} - \frac{1}{2}\eta_{\mu\nu}h$.

The general solution includes all six polarizations:
\begin{equation}
    h_{\mu\nu}(t,\mathbf{x}) = \sum_{A=1}^6 \int \frac{d^3k}{(2\pi)^3} \left[ h_A(\mathbf{k}) e^{i(k \cdot x - \omega_A t)} + h_A^*(\mathbf{k}) e^{-i(k \cdot x - \omega_A t)} \right] e^A_{\mu\nu}(\hat{k})
\end{equation}

where $\omega_A = v_A |\mathbf{k}|$ depends on the polarization.

\subsection{Strong Field Regime}

In the strong field regime (near black holes), torsion effects become significant:
\begin{equation}
    \frac{|T|}{|R|} \sim \tau_0 \left(\frac{M_P}{M}\right)^{1/2}
\end{equation}

For stellar mass black holes ($M \sim 10 M_\odot$):
\begin{equation}
    \frac{|T|}{|R|} \sim 10^{-6} \times 10^{-19} \sim 10^{-25}
\end{equation}

This is extremely small, but for primordial black holes ($M \sim 10^{-5}$ g):
\begin{equation}
    \frac{|T|}{|R|} \sim 10^{-6} \times 10^{19} \sim 10^{13}
\end{equation}

Torsion dominates for microscopic black holes.

\subsection{Detection Prospects}

The six polarization modes can be detected through:

\textbf{1. LISA (Space-based)}:
- Frequency range: $10^{-4}$ Hz to 1 Hz
- Sensitive to intermediate mass black hole mergers ($10^4$-$10^7 M_\odot$)
- Can resolve vector polarizations with SNR $\sim$ 7 for $\tau_0 = 10^{-6}$

\textbf{2. Cosmic Explorer/Einstein Telescope}:
- Frequency range: 1 Hz to 10 kHz
- Ground-based, higher frequencies
- Complementary to LISA

\textbf{3. Pulsar Timing Arrays}:
- Frequency range: nHz to $\mu$Hz
- Stochastic background detection
- Sensitive to supermassive black hole binaries

\subsection{Correspondence with General Relativity}

Our theory smoothly reduces to general relativity when:
\begin{enumerate}
    \item $\tau_0 \to 0$: Torsion effects vanish
    \item $E \ll E_c$: Internal space decouples
    \item Weak field: Non-linear torsion terms negligible
\end{enumerate}

All classical tests of GR (perihelion precession, light bending, gravitational redshift, Shapiro delay) are satisfied in the appropriate limits.

\subsection{Summary}

The gravitational theory in our framework:
\begin{itemize}
    \item Extends GR to Einstein-Cartan theory with torsion
    \item Predicts six gravitational wave polarizations
    \item Vector polarizations have amplitude $\sim \tau_0/2$
    \item Scalar polarizations have amplitude $\sim \tau_0^2/3$
    \item Different propagation speeds lead to detectable time delays
    \item Reduces to GR in the low-energy, weak-field limit
\end{itemize}

The next section derives the gauge theory of electromagnetism, weak and strong nuclear forces from the same geometric framework.
